\Needspace{12\baselineskip}
\begin{center}
\textit{Special applications to writing numbers as sums of two squares, a square and twice a square, a square and three times a square.}
\end{center}

\Needspace{8\baselineskip}
\Article{182}
We turn to some special cases, both for their remarkable elegance and for the sustained attention the distinguished Euler gave them, which has earned them something of the status of classics.

I. A number can be represented by the form $xx+yy$ with $x$ relatively prime to $y$ (or written as a sum of
\SourceLeaf{165}
two relatively prime squares) only if $-1$ is a quadratic residue of it. Conversely, every such number, taken positively, can be represented. Let $M$ be such a number, and let all values of the expression $\surd-1\ (\mathrm{mod.}\ M)$ be $N,-N,N',-N',N'',-N''$ etc. Then, by Article 176, the form $(M,N,\frac{NN+1}{M})$ is properly equivalent to the form $(1,0,1)$. Let a proper transformation of the latter into the former be $x=\alpha x'+\beta y'$, $y=\gamma x'+\delta y'$. The representations of the number $M$ by the form $xx+yy$ belonging to $N$ are then these four\HistoricalNote{*)}{This case is clearly covered by (2) of Article {180}.}: $x=\pm\alpha,y=\pm\gamma$; $x=\mp\gamma,y=\pm\alpha$.

Since the form $(1,0,1)$ is ambiguous, the form $(M,-N,\frac{NN+1}{M})$ is clearly also properly equivalent to it, and the former changes properly into the latter on putting $x=\alpha x'-\beta y'$, $y=-\gamma x'+\delta y'$. This gives four representations of $M$ belonging to $-N$: $x=\pm\alpha,y=\mp\gamma$; $x=\pm\gamma,y=\pm\alpha$. So there are eight representations of $M$, half belonging to $N$ and half to $-N$. But all these give only \textit{one} decomposition of the number $M$ into two squares, $M=\alpha\alpha+\gamma\gamma$, if we consider only the squares themselves, without regard to their order or the signs of their roots.

So if there are no values of the expression $\surd-1\ (\mathrm{mod.}\ M)$ other than $N$ and $-N$, as happens, \textit{for example}, when $M$ is prime, $M$ can be written as a sum of two relatively prime squares in only one way. Now $-1$ is a quadratic residue of every prime of the form $4n+1$ (Article 108), and a prime clearly cannot be written as a sum of two squares that are not relatively prime. We therefore have the theorem:

\textit{Every prime of the form $4n+1$ can be written as a sum of two squares, in only one way.}
\[
\begin{gathered}
1=0+1,\quad5=1+4,\quad13=4+9,\quad17=1+16,\\
29=4+25,\quad37=1+36,\quad41=16+25,\quad53=4+49,\\
61=25+36,\quad73=9+64,\quad89=25+64,\quad97=16+81\quad\text{etc.}
\end{gathered}
\]
This very elegant theorem was already known to Fermat, but was first proved by the distinguished Euler, \textit{Comm. nov. Petr. Vol.}\ V, \textit{for the years} 1754, 1755, \textit{p.}\ 3 ff. \textit{Vol.}\ IV contains a paper on the same subject, \textit{p.}\ 3 ff., but he had not yet completed the matter there; see especially Article 27.

\SourceLeaf{166}
So if a number of the form $4n+1$ can be written as a sum of two squares either in several ways or in no way, it is certainly not prime.

\Needspace{5\baselineskip}
Conversely, if the expression $\surd-1\ (\mathrm{mod.}\ M)$ has values other than $N$ and $-N$, there are further representations of $M$ belonging to them. In this case, then, $M$ can be written as a sum of two squares in several ways, \textit{for example}, $65=1+64=16+49$, $221=25+196=100+121$.

The remaining representations, in which $x,y$ take values that are not relatively prime, can easily be found by our general method. We only note that if a number has factors of the form $4n+3$ that cannot be removed by dividing by any square (which happens if one or more such factors has an \textit{odd exponent}), it cannot be written as a sum of two squares in any way\HistoricalNote{*)}{If the number $M=2^{\mu}Sa^{\alpha}b^{\beta}c^{{\gamma}}\ldots$, where $a,b,c$ etc. are distinct primes of the form $4n+1$, and $S$ is the product of all prime factors of $M$ of the form $4n+3$ (every positive number can be put in this form by taking $\mu=0$ when $M$ is odd, and $S=1$ when $M$ has no factors of the form $4n+3$), then $M$ can be written as a sum of two squares in \textit{no way} if $S$ is not a square. If $S$ is a square, there are $\tfrac12(\alpha+1)(\beta+1)(\gamma+1)$ etc. decompositions of $M$ when at least one of the numbers $\alpha,\beta,\gamma$ etc. is odd, or $\tfrac12(\alpha+1)(\beta+1)(\gamma+1)\text{ etc.}+\tfrac12$ when all $\alpha,\beta,\gamma$ etc. are even (considering only the squares themselves). Anyone with some practice in combinatorics can derive the proof of this theorem from our general theory without much difficulty. We cannot dwell on it here, any more than on other \textit{special cases}. Compare Article 105.}.

II. The form $xx+2yy$ cannot represent a number for which $-2$ is a nonresidue with $x$ relatively prime to $y$; every other number can be represented. Let $-2$ be a residue of the number $M$, and $N$ a value of the expression $\surd-2\ (\mathrm{mod.}\ M)$. Then, by Article 176, the forms $(1,0,2)$, $(M,N,\frac{NN+2}{M})$ are properly equivalent. Let the first change properly into the second on putting $x=\alpha x'+\beta y'$, $y=\gamma x'+\delta y'$. Then $x=\alpha,y=\gamma$ is a representation of the number $M$ belonging to $N$. Except this one and $x=-\alpha,y=-\gamma$, no others belong to $N$ (Article 180).

As above, we see that the representations $x=\pm\alpha,y=\mp\gamma$ belong to the value $-N$. All four representations give just one decomposition of $M$ into a square and twice a square. If there are no values other than $N$ and $-N$ of the expression $\surd-2\ (\mathrm{mod.}\ M)$, there are no other decompositions. With the propositions of Article 116, this easily gives the theorem:

\SourceLeaf{167}
\textit{Every prime of the form $8n+1$ or $8n+3$ can be written as a square plus twice a square, in only one way.}
\[
\begin{gathered}
1=1+0,\quad3=1+2,\quad11=9+2,\quad17=9+8,\\
19=1+18,\quad41=9+32.\quad43=25+18,\quad59=9+50,\\
67=49+18,\quad73=1+72,\quad83=81+2,\\
89=81+8,\quad97=25+72\quad\text{etc.}
\end{gathered}
\]
Fermat also knew this theorem, like several similar ones, but the distinguished La Grange first gave a proof, \textit{Suite des recherches d'Arithmétique}, Nouv. Mém. de l'Ac. de Berlin 1775, p. 323 ff. The distinguished Euler had already settled many points on this subject, \textit{Specimen de usu observationum in mathesi pura}, Comm. nov. Petr. Vol. VI, p. 185 ff. But a complete proof of the theorem always eluded his efforts, p. 220. Compare also the paper in Vol. VIII (for the years 1760, 1761), \textit{Supplementum quorundam theorematum arithmeticorum}, near the end.

\Needspace{12\baselineskip}
III. A similar method proves that every number of which $-3$ is a quadratic residue can be represented either by the form $xx+3yy$ or by $2xx+2xy+2yy$, with the value of $x$ relatively prime to the value of $y$. Since $-3$ is a residue of all primes of the form $3n+1$ (Article 119), and the form $2xx+2xy+2yy$ clearly represents only \textit{even} numbers, we have, as above, the theorem:

\textit{Every prime of the form $3n+1$ can be written as a square plus three times a square, in only one way.}
\[
\begin{gathered}
1=1+0,\quad7=4+3,\quad13=1+12,\quad19=16+3,\\
31=4+27,\quad37=25+12,\quad43=16+27,\quad61=49+12,\\
67=64+3,\quad73=25+48\quad\text{etc.}
\end{gathered}
\]
The distinguished Euler first gave a proof of this theorem in the paper just cited, \textit{Comm. nov. Petr. Vol.}\ VIII, \textit{p.}\ 105 ff.

We could proceed further in the same way and, \textit{for example}, show that every prime of the form $20n+1$, $20n+3$, $20n+7$, or $20n+9$ (of which $-5$ is a residue) can be represented by one of the forms $xx+5yy$, $2xx+2xy+3yy$: primes of the forms $20n+1$ and $20n+9$ by the first, primes of the forms $20n+3,20n+7$ by the second. Also, twice a prime of the forms $20n+1,20n+9$ can be represented by the form $2xx+2xy+3yy$, and twice a prime of the forms $20n+3,20n+7$ by the form $xx+5yy$. But anyone can derive this proposition and infinitely many other special ones for themselves from the preceding results and those to be given
\SourceLeaf{168}
below. --- So we turn to \textit{forms with positive determinant}. Since their nature is quite different according as the determinant is a square or not a square, we first exclude forms with square determinant here and will consider them separately later.
